\documentclass[10pt,a4paper]{amsart}
\usepackage[margin=19mm]{geometry}
\usepackage{amsmath,amssymb,mathtools}
\usepackage[scr=rsfs,cal=euler]{mathalfa}
\usepackage{xfrac}
\usepackage[unicode,hidelinks,bookmarks=false]{hyperref}
\newcommand{\ie}{i.e.\ }
\newcommand{\cf}{cf.\ }
\newcommand{\resp}{respectively\ }
\newcommand{\Subsection}{\subsection}
\providecommand{\nobreakdash}{\nobreak\hbox{-}}
\setlength{\emergencystretch}{2em}
\multlinegap0pt
\usepackage{tikz}
\usetikzlibrary{cd}
\usepackage{amssymb}
\usepackage[shortlabels]{enumitem}
\newtheorem{prop}{Proposition}
\newcommand{\dcat}{\mathcal{D}}
\newcommand{\ftopos}{\mathcal{F}}
\newcommand{\locrefl}{\mathcal{L}}
\newcommand{\pretopos}{\mathcal{P}}
\newcommand{\topos}{\mathcal{E}}
\title{A counterexample to Kanalas' problem of continuously realising types}
\author{Morgan Rogers, Joshua Wrigley}
\date{Source: arXiv:2608.25688v1 (CC BY 4.0)}
\begin{document}
\maketitle

\noindent\emph{Source and licence.} A counterexample to Kanalas' problem of continuously realising types. Version arXiv:2608.25688v1 (CC BY 4.0), distributed by arXiv under the Creative Commons Attribution 4.0 International licence, http://creativecommons.org/licenses/by/4.0/, as recorded in the arXiv licence metadata for this exact version and shown on the abstract page https://arxiv.org/abs/2608.25688v1. Full licence notice: \texttt{licenses/arxiv-2608.25688-CC-BY-4.0.txt}.\par\smallskip

\noindent\textit{Editorial scope.} Counted unit: the article's prop prop (source0.tex lines 142--152). The article's complete original proof of it is delivered with it (source0.tex lines 153--193, one \texttt{proof} environment of 41 lines). No mathematics is written, repaired or re-derived: the statement and the proof are the article's own lines, in the article's own notation, and no retained line is substituted or re-flowed.  No further source-local context is needed: every cross-reference of every retained line resolves inside the two delivered spans.  Nothing was omitted from any retained span, so the package carries no omission record.  No retained line contains a citation, so the wrapper carries no bibliography and no bibliography entry.  The retained lines contain the article's own diagram code, which the wrapper typesets with the standard tikz packages; no image file is delivered and the package declares no asset dependency.  Editorial formatting only, in addition: an AMS \texttt{amsart} preamble that restates the article's own declarations and macros used by the retained lines, namely the environments \texttt{prop} on separate counters, together with the standard packages those lines need.  The retained environments print in this document as Proposition 1 in order of appearance, whereas the article prints the same items with the numbering of its own full document.  The complete native source of the article as distributed in the arXiv e-print for this exact version is bundled unchanged as \texttt{sources/208589/source0.tex}.
	\begin{prop}
		Let $\pretopos$ be a pretopos with classifying topos $\topos$.  Every Grothendieck topos realises types for $\pretopos$ if and only if the localic reflection $h \colon \topos \to \locrefl$ admits a lax section, as in the following diagram.
		\[\begin{tikzcd}
			\locrefl & \topos \\
			& \locrefl
			\arrow[dashed, from=1-1, to=1-2]
			\arrow[""{name=0, anchor=center, inner sep=0}, equals, from=1-1, to=2-2]
			\arrow["h", from=1-2, to=2-2]
			\arrow["\Rightarrow"{marking, allow upside down}, draw=none, from=0, to=1-2]
		\end{tikzcd}\]
	\end{prop}

	\begin{proof}
		Recall that $\dcat$ embeds into its classifying topos $\locrefl$, and recall also that this inclusion $\dcat \hookrightarrow \locrefl$ corresponds, under the classifying property of $\locrefl$, to the identity morphism on $\locrefl$.  If $\locrefl$ realises types, then there must be a natural transformation
		\[\begin{tikzcd}
			\dcat & \\
			\pretopos & \locrefl,
			\arrow[hook, from=1-1, to=2-1]
			\arrow[""{name=0, anchor=center, inner sep=0}, hook, from=1-1, to=2-2]
			\arrow[dashed, from=2-1, to=2-2]
			\arrow["\Rightarrow"{marking, allow upside down}, draw=none, from=0, to=2-1]
		\end{tikzcd}\]
		which, via converting to classifying topoi, yields the desired lax splitting
		\[\begin{tikzcd}
			\locrefl & \\
			\topos & \locrefl.
			\arrow[""{name=0, anchor=center, inner sep=0}, equals, from=1-1, to=2-2]
			\arrow["h", from=2-1, to=1-1]
			\arrow[dashed, from=2-2, to=2-1]
			\arrow["\Rightarrow"{marking, allow upside down}, draw=none, from=0, to=2-1]
		\end{tikzcd}\]
		
		Now suppose we are given a lax section as above.  Then for any other topos $\ftopos$ and a coherent functor $\dcat \to \ftopos$, we take the corresponding geometric morphism $\ftopos \to \locrefl$ and then form the following pasting of geometric morphisms.
		\[\begin{tikzcd}
			\ftopos & \locrefl & \topos \\
			&& \locrefl
			\arrow[from=1-1, to=1-2]
			\arrow[from=1-2, to=1-3]
			\arrow[""{name=0, anchor=center, inner sep=0}, equals, from=1-2, to=2-3]
			\arrow["h", from=1-3, to=2-3]
			\arrow["\Rightarrow"{marking, allow upside down}, draw=none, from=0, to=1-3]
		\end{tikzcd}\]
		Now, converting back to the coherent functors that are being classified, we obtain a diagram of coherent functors
		\[\begin{tikzcd}
			\ftopos & \pretopos \\
			& \dcat
			\arrow[from=1-2, to=1-1]
			\arrow[""{name=0, anchor=center, inner sep=0}, from=2-2, to=1-1]
			\arrow[hook, from=2-2, to=1-2]
			\arrow["\Rightarrow"{marking, allow upside down}, draw=none, from=0, to=1-2]
		\end{tikzcd}\]
		as desired.
	\end{proof}

\end{document}
